%! AP Statistics — Unit 3, Lesson 3.5 — Homework
\documentclass[10pt]{article}
\usepackage{apstats-article}
\usepackage{apstats-boxes}

\begin{document}

\pageheader{Unit 3, Lesson 3.5}{Homework}
\namedateperiod

\vspace{0.1in}

\begin{practicebox}
\small
For each experiment, identify units, factor(s), treatments, response variable, design
type, and blinding used.

\vspace{8pt}
\textbf{1.} A study tests whether a new sleep aid improves sleep quality. 100 adults
are randomly assigned: 50 take the pill nightly; 50 take a look-alike sugar pill.
Participants don't know which pill they take. A sleep technician, also unaware of
treatment, rates sleep quality.

\vspace{4pt}
Experimental units: \blank{8cm}\\[4pt]
Factor: \blank{8cm}\\[4pt]
Treatments: \blank{8cm}\\[4pt]
Response variable: \blank{8cm}\\[4pt]
Design type: \blank{4cm} \quad Blinding: \blank{4cm}\\[4pt]
Why is the placebo important? \writeline\\[1pt]\writeline

\vspace{8pt}
\textbf{2.} A farmer wants to compare 3 fertilizers. He has 90 plants all grown in
the same greenhouse. He randomly assigns 30 plants to each fertilizer.

\vspace{4pt}
Design type: \blank{4cm} \quad Blocking variable (if any): \blank{4cm}\\[4pt]
Would adding blocks improve this design? Explain. \writeline\\[1pt]\writeline

\vspace{8pt}
\textbf{3.} A researcher tests two diets on 60 participants. She pairs participants
by current weight (closest pairs); one in each pair gets Diet A, one gets Diet B
(randomly assigned).

\vspace{4pt}
Design type: \blank{7cm}\\[4pt]
Why is pairing by weight appropriate here? \writeline\\[1pt]\writeline

\vspace{8pt}
\textbf{4.} Reflection: Why does random \emph{assignment} (not just random selection)
allow us to conclude causation from an experiment?\\[4pt]
\writeline\\[1pt]\writeline\\[1pt]\writeline
\end{practicebox}

\end{document}
